\documentclass[UTF8]{ctexart}
\usepackage{xeCJK}
\usepackage{geometry}
\geometry{left=2cm,right=2cm,top=2.5cm,bottom=2.5cm}
\setlength{\headheight}{12.7pt}

\usepackage{titlesec}
\titleformat{\section}
  {\normalfont\large\bfseries}
  {\thesection}
  {1em}
  {}
\titlespacing*{\section}
  {0pt}
  {3.5ex plus 1ex minus .2ex}
  {2.3ex plus .2ex}

\usepackage{amsmath}
\usepackage{amssymb}
\usepackage{amsthm}
\usepackage{enumitem}
\setlist[enumerate]{itemsep=0pt,topsep=0pt,parsep=0pt,leftmargin=0pt,itemindent=2.5em}
\setlist[itemize]{itemsep=0pt,topsep=0pt,parsep=0pt,leftmargin=0pt,itemindent=2.5em}
\usepackage{fancyhdr}
\pagestyle{fancy}
\fancyhf{}
\usepackage{graphicx}
\usepackage{hyperref}
\usepackage{indentfirst}
\usepackage{lmodern}


\begin{document}
\setlength{\abovedisplayskip}{1pt}
\setlength{\belowdisplayskip}{1pt}

\section{while 程序操作语义}

\hypertarget{sec:定义1.1}{}
\noindent\textbf{定义1.1 (配置)}

固定一个模型$M$, 记$\Sigma$为所有变量赋值(状态)的集合, 定义
\[
    (\,\mathcal{C}\cup\{\mathrm{E}\}\,)\times\Sigma
\]
为\textbf{配置集}, 其中$\langle S,\sigma\rangle$($S\in\mathcal{C}$)
称作未终止的配置, $\langle\mathrm{E},\sigma\rangle$称作已终止的配置.

\vspace{10pt}

\hypertarget{sec:定义1.2}{}
\noindent\textbf{定义1.2 (一步转移操作语义)}

递归定义未终止的配置上的一步转移关系
$\mathcal{C}\times\Sigma\to(\,\mathcal{C}\cup\{\mathrm{E}\}\,)\times\Sigma$ :
\begin{enumerate}
    \item 对任意的状态$\sigma$,
    \vspace{-3pt}
    \[
        \langle\mathrm{skip},\sigma\rangle\overset{\mathrm{def}}{\to}
        \langle\mathrm{E},\sigma\rangle,\\[-3pt]
    \]
    \item 对任意同类型的变元$u$, 项$t$和状态$\sigma$,
    \[
        \langle u:=t,\sigma\rangle\overset{\mathrm{def}}{\to}
        \langle\mathrm{E},\sigma[u:=\sigma(t)]\rangle,\\[3pt]
    \]
    \item 对任意的$S_1,S_2$和状态$\sigma$,
    记$\langle S_1,\sigma\rangle\to\langle S_1',\sigma'\rangle$, 则
    \vspace{3pt}
    \[
        \langle S_1;S_2,\sigma\rangle\overset{\mathrm{def}}{\to}
        \langle S_1';S_2,\sigma'\rangle,\\[3pt]
    \]
    其中$\mathrm{E};S_2$定义为$S_2$.
    \item 对任意的$B:\mathrm{Bool}$和$S_1,S_2\in\mathcal{C}$,
    如果$\sigma\models B$, 则
    \[
        \langle\mathbf{if}\,\,B\,\,\mathbf{then}\,\,S_1\,\,\mathbf{else}\,\,S_2
        \,\,\mathbf{fi},\sigma\rangle\overset{\mathrm{def}}{\to}
        \langle S_1,\sigma\rangle,
    \]
    \item 对任意的$B:\mathrm{Bool}$和$S_1,S_2\in\mathcal{C}$,
    如果$\sigma\not\models B$, 则
    \[
        \langle\mathbf{if}\,\,B\,\,\mathbf{then}\,\,S_1\,\,\mathbf{else}\,\,S_2
        \,\,\mathbf{fi},\sigma\rangle\overset{\mathrm{def}}{\to}
        \langle S_2,\sigma\rangle,
    \]
    \item 对任意的$B:\mathrm{Bool}$和$S\in\mathcal{C}$,
    如果$\sigma\models B$, 则
    \[
        \langle\mathbf{while}\,\,B\,\,\mathbf{do}\,\,S\,\,\mathbf{od},\sigma\rangle
        \overset{\mathrm{def}}{\to}\langle S;
        \,\mathbf{while}\,\,B\,\,\mathbf{do}\,\,S\,\,\mathbf{od},\sigma\rangle,
    \]
    \item 对任意的$B:\mathrm{Bool}$和$S\in\mathcal{C}$,
    如果$\sigma\not\models B$, 则
    \[
        \langle\mathbf{while}\,\,B\,\,\mathbf{do}\,\,S\,\,\mathbf{od},\sigma\rangle
        \overset{\mathrm{def}}{\to}\langle\mathrm{E},\sigma\rangle.
    \]
\end{enumerate}





\newpage

\hypertarget{sec:定义1.3}{}
\noindent\textbf{定义1.3 (转移序列与计算)}

\begin{enumerate}
    \item 给定有限的配置序列$\langle S_i,\sigma_i\rangle$ ($i<n,\,n\geqslant 1$),
    如果对每个$i<n-1$有
    \[
        \langle S_i,\sigma_i\rangle\to\langle S_{i+1},\sigma_{i+1}\rangle,
    \]
    就称它是从$\langle S_0,\sigma_0\rangle$开始的一个\textbf{有限转移序列}.

    \hspace{2.17em}
    当$S_{n-1}=\mathrm{E}$时也称作一个\textbf{计算},
    称其终止于$\sigma_{n-1}$.

    \item 给定无限的配置序列$\langle S_i,\sigma_i\rangle$ ($i\in\mathbb{N}$),
    如果对每个$i\in\mathbb{N}$有
    \[
        \langle S_i,\sigma_i\rangle\to\langle S_{i+1},\sigma_{i+1}\rangle,
    \]
    就称它是从$\langle S_0,\sigma_0\rangle$开始的一个\textbf{无限转移序列},
    也称作一个发散的\textbf{计算}.
\end{enumerate}

\vspace{10pt}

\hypertarget{sec:定理1.1}{}
\noindent\textbf{定理1.1}

任给程序$S$和状态$\sigma$, 存在唯一的从$\langle S,\sigma\rangle$开始的计算.

\noindent{\kaishu 证明.}

\noindent{\kaishu 存在性.}

对$i\in\mathbb{N}$递归定义$S_i,\sigma_i$:
\[\langle S_i,\sigma_i\rangle=\left\{\begin{aligned}
    &\langle S,\sigma\rangle,&&i=0,\\
    &\to\langle S_{i-1},\sigma_{i-1}\rangle,&&i>0\,\land\,S_{i-1}\in\mathcal{C},\\
    &\langle\mathrm{E},\sigma_{i-1}\rangle,&&i>0\,\land\,S_{i-1}=\mathrm{E}.
\end{aligned}\right.\]

\vspace{3pt}
如果序列中出现了$\mathrm{E}$, 则取最小的使得$S_n=\mathrm{E}$的$n$,
易证$\langle S_i,\sigma_i\rangle$ ($i\leqslant n$)是有限的计算;

如果序列中没有出现$\mathrm{E}$, 则易证整个序列是无限的计算.

\noindent{\kaishu 唯一性.}

归纳证明转移序列的每一位都确定即可. \hfill $\qedsymbol$

\end{document}
